\documentclass[a4paper,12pt,twoside]{memoir}

% Castellano
\usepackage[spanish,es-tabla]{babel}
\selectlanguage{spanish}
\usepackage[utf8]{inputenc}
\usepackage[T1]{fontenc}
\usepackage{lmodern} % scalable font
\usepackage{microtype}
\usepackage{placeins}

\RequirePackage{booktabs}
\RequirePackage[table]{xcolor}
\RequirePackage{xtab}
\RequirePackage{multirow}

% Links
\PassOptionsToPackage{hyphens}{url}\usepackage[colorlinks]{hyperref}
\hypersetup{
	allcolors = {red}
}

% Ecuaciones
\usepackage{amsmath}

% Rutas de fichero / paquete
\newcommand{\ruta}[1]{{\sffamily #1}}

% Párrafos
\nonzeroparskip

% Huérfanas y viudas
\widowpenalty100000
\clubpenalty100000

% Evitar solapes en el header
\nouppercaseheads

% Imagenes
\usepackage{graphicx}
\newcommand{\imagen}[2]{
	\begin{figure}[!h]
		\centering
		\includegraphics[width=0.9\textwidth]{#1}
		\caption{#2}\label{fig:#1}
	\end{figure}
	\FloatBarrier
}

\newcommand{\imagenflotante}[2]{
	\begin{figure}%[!h]
		\centering
		\includegraphics[width=0.9\textwidth]{#1}
		\caption{#2}\label{fig:#1}
	\end{figure}
}



% El comando \figura nos permite insertar figuras comodamente, y utilizando
% siempre el mismo formato. Los parametros son:
% 1 -> Porcentaje del ancho de página que ocupará la figura (de 0 a 1)
% 2 --> Fichero de la imagen
% 3 --> Texto a pie de imagen
% 4 --> Etiqueta (label) para referencias
% 5 --> Opciones que queramos pasarle al \includegraphics
% 6 --> Opciones de posicionamiento a pasarle a \begin{figure}
\newcommand{\figuraConPosicion}[6]{%
  \setlength{\anchoFloat}{#1\textwidth}%
  \addtolength{\anchoFloat}{-4\fboxsep}%
  \setlength{\anchoFigura}{\anchoFloat}%
  \begin{figure}[#6]
    \begin{center}%
      \Ovalbox{%
        \begin{minipage}{\anchoFloat}%
          \begin{center}%
            \includegraphics[width=\anchoFigura,#5]{#2}%
            \caption{#3}%
            \label{#4}%
          \end{center}%
        \end{minipage}
      }%
    \end{center}%
  \end{figure}%
}

%
% Comando para incluir imágenes en formato apaisado (sin marco).
\newcommand{\figuraApaisadaSinMarco}[5]{%
  \begin{figure}%
    \begin{center}%
    \includegraphics[angle=90,height=#1\textheight,#5]{#2}%
    \caption{#3}%
    \label{#4}%
    \end{center}%
  \end{figure}%
}
% Para las tablas
\newcommand{\otoprule}{\midrule [\heavyrulewidth]}
%
% Nuevo comando para tablas pequeñas (menos de una página).
\newcommand{\tablaSmall}[5]{%
 \begin{table}
  \begin{center}
   \rowcolors {2}{gray!35}{}
   \begin{tabular}{#2}
    \toprule
    #4
    \otoprule
    #5
    \bottomrule
   \end{tabular}
   \caption{#1}
   \label{tabla:#3}
  \end{center}
 \end{table}
}

%
%Para el float H de tablaSmallSinColores
\usepackage{float}

%
% Nuevo comando para tablas pequeñas (menos de una página).
\newcommand{\tablaSmallSinColores}[5]{%
 \begin{table}[H]
  \begin{center}
   \begin{tabular}{#2}
    \toprule
    #4
    \otoprule
    #5
    \bottomrule
   \end{tabular}
   \caption{#1}
   \label{tabla:#3}
  \end{center}
 \end{table}
}

\newcommand{\tablaApaisadaSmall}[5]{%
\begin{landscape}
  \begin{table}
   \begin{center}
    \rowcolors {2}{gray!35}{}
    \begin{tabular}{#2}
     \toprule
     #4
     \otoprule
     #5
     \bottomrule
    \end{tabular}
    \caption{#1}
    \label{tabla:#3}
   \end{center}
  \end{table}
\end{landscape}
}

%
% Nuevo comando para tablas grandes con cabecera y filas alternas coloreadas en gris.
\newcommand{\tabla}[6]{%
  \begin{center}
    \tablefirsthead{
      \toprule
      #5
      \otoprule
    }
    \tablehead{
      \multicolumn{#3}{l}{\small\sl continúa desde la página anterior}\\
      \toprule
      #5
      \otoprule
    }
    \tabletail{
      \hline
      \multicolumn{#3}{r}{\small\sl continúa en la página siguiente}\\
    }
    \tablelasttail{
      \hline
    }
    \bottomcaption{#1}
    \rowcolors {2}{gray!35}{}
    \begin{xtabular}{#2}
      #6
      \bottomrule
    \end{xtabular}
    \label{tabla:#4}
  \end{center}
}

%
% Nuevo comando para tablas grandes con cabecera.
\newcommand{\tablaSinColores}[6]{%
  \begin{center}
    \tablefirsthead{
      \toprule
      #5
      \otoprule
    }
    \tablehead{
      \multicolumn{#3}{l}{\small\sl continúa desde la página anterior}\\
      \toprule
      #5
      \otoprule
    }
    \tabletail{
      \hline
      \multicolumn{#3}{r}{\small\sl continúa en la página siguiente}\\
    }
    \tablelasttail{
      \hline
    }
    \bottomcaption{#1}
    \begin{xtabular}{#2}
      #6
      \bottomrule
    \end{xtabular}
    \label{tabla:#4}
  \end{center}
}

%
% Nuevo comando para tablas grandes sin cabecera.
\newcommand{\tablaSinCabecera}[5]{%
  \begin{center}
    \tablefirsthead{
      \toprule
    }
    \tablehead{
      \multicolumn{#3}{l}{\small\sl continúa desde la página anterior}\\
      \hline
    }
    \tabletail{
      \hline
      \multicolumn{#3}{r}{\small\sl continúa en la página siguiente}\\
    }
    \tablelasttail{
      \hline
    }
    \bottomcaption{#1}
  \begin{xtabular}{#2}
    #5
   \bottomrule
  \end{xtabular}
  \label{tabla:#4}
  \end{center}
}



\definecolor{cgoLight}{HTML}{EEEEEE}
\definecolor{cgoExtralight}{HTML}{FFFFFF}

%
% Nuevo comando para tablas grandes sin cabecera.
\newcommand{\tablaSinCabeceraConBandas}[5]{%
  \begin{center}
    \tablefirsthead{
      \toprule
    }
    \tablehead{
      \multicolumn{#3}{l}{\small\sl continúa desde la página anterior}\\
      \hline
    }
    \tabletail{
      \hline
      \multicolumn{#3}{r}{\small\sl continúa en la página siguiente}\\
    }
    \tablelasttail{
      \hline
    }
    \bottomcaption{#1}
    \rowcolors[]{1}{cgoExtralight}{cgoLight}

  \begin{xtabular}{#2}
    #5
   \bottomrule
  \end{xtabular}
  \label{tabla:#4}
  \end{center}
}




\graphicspath{ {../imagenes/} }

% Capítulos
\chapterstyle{bianchi}
\newcommand{\capitulo}[2]{
	\setcounter{chapter}{#1}
	\setcounter{section}{0}
	\setcounter{figure}{0}
	\setcounter{table}{0}
	\chapter*{#2}
	\addcontentsline{toc}{chapter}{#2}
	\markboth{#2}{#2}
}

% Apéndices
\renewcommand{\appendixname}{Apéndice}
\renewcommand*\cftappendixname{\appendixname}

\newcommand{\apendice}[1]{
	%\renewcommand{\thechapter}{A}
	\chapter{#1}
}

\renewcommand*\cftappendixname{\appendixname\ }

% Formato de portada
\makeatletter
\usepackage{xcolor}
\newcommand{\tutor}[1]{\def\@tutor{#1}}
\newcommand{\course}[1]{\def\@course{#1}}
\definecolor{cpardoBox}{HTML}{E6E6FF}
\def\maketitle{
  \null
  \thispagestyle{empty}
  % Cabecera ----------------
\noindent\includegraphics[width=\textwidth]{cabecera}\vspace{1cm}%
  \vfill
  % Título proyecto y escudo informática ----------------
  \colorbox{cpardoBox}{%
    \begin{minipage}{.8\textwidth}
      \vspace{.5cm}\Large
      \begin{center}
      \textbf{TFG del Grado en Ingeniería Informática}\vspace{.6cm}\\
      \textbf{\LARGE\@title{}}
      \end{center}
      \vspace{.2cm}
    \end{minipage}

  }%
  \hfill\begin{minipage}{.20\textwidth}
    \includegraphics[width=\textwidth]{escudoInfor}
  \end{minipage}
  \vfill
  % Datos de alumno, curso y tutores ------------------
  \begin{center}%
  {%
    \noindent\LARGE
    Presentado por \@author{}\\ 
    en Universidad de Burgos --- \@date{}\\
    Tutor: \@tutor{}\\
  }%
  \end{center}%
  \null
  \cleardoublepage
  }
\makeatother


% Datos de portada
\title{título del TFG \\Documentación Técnica}
\author{José Luis Pérez Gómez}
\tutor{Bruno Baruque}
\date{\today}

\begin{document}

\maketitle



\cleardoublepage



%%%%%%%%%%%%%%%%%%%%%%%%%%%%%%%%%%%%%%%%%%%%%%%%%%%%%%%%%%%%%%%%%%%%%%%%%%%%%%%%%%%%%%%%



\frontmatter


\clearpage

% Indices
\tableofcontents

\clearpage

\listoffigures

\clearpage

\listoftables

\clearpage

\mainmatter

\appendix

\apendice{Plan de Proyecto Software}

\section{Introducción}

La Universidad de Burgos, dentro del área de conocimiento de Ingeniería de Sistemas y Automática, dispone de un interfaz BCI (Brain Computer Interface) para la captación de señales cerebrales. 
Empleando ese interfaz se han realizado diferentes experimentos que han permitido recoger información de la actividad cerebral mientras los usuarios ejecutaban diferentes tareas cotidianas. 

Este Trabajo de Fin de Grado (TFG) tiene como objetivo principal el análisis de la información obtenida en esos experimentos. Se entrenarán diferentes algoritmos para clasificar la acción realizada por el usuario a partir de las señales generadas por el BCI. Con este propósito, se evaluarán diferentes algoritmos de procesado de señales y de machine/deep learning para la clasificación automática de señales.

El conjunto de datos para la realización del TFG son de tipo EEG (Electroencefalografía) referentes a experimentos basados en acciones sobre teclas de un teclado: arriba, abajo, izquierda, derecha.

El análisis de estos datos y su evaluación en diferentes algoritmos está basada en predecir qué teclas del teclado se han pulsado según las señales captadas con la interfaz BCI.



\section{Planificación temporal}


En la reunión inicial con los tutores definimos utilizar Python como lenguaje de programación para escribir el código y una serie de cursos de aprendizaje básicos para poder acometer el TFG sin problemas.

En las siguientes reuniones se definieron algoritmos y experimentos a realizar con los datos EEG recogidos.


Desde el inicio del proyecto se ha utilizado la herramienta GitHub para ir subiendo todo el código o documentación realizada vía commits al repositorio creado para el proyecto ~\cite{github:proyecto}, a continuación se muestra la linea temporal durante la elaboración del proyecto:

\imagen{anexos/Commits1}{Commits realizados durante la realización del TFG}

\imagen{anexos/Additions1}{Additions realizados durante la realización del TFG}

\imagen{anexos/Deletions1.png}{Deletions realizados durante la realización del TFG}


El desarrollo de este proyecto siguió la metodología CRISP-DM (Cross Industry Standard Process for Data Mining) ~\cite{wiki:CRISP}, un enfoque estructurado y bien establecido para proyectos de minería de datos y análisis predictivo. 

A continuación, se detallan los aspectos más interesantes y relevantes de cada fase del ciclo de vida del proyecto según CRISP-DM.


\begin{enumerate}
\def\labelenumi{\arabic{enumi}.}	
\item
\textbf{Comprensión del Negocio}

El objetivo principal fue identificar los objetivos y requisitos del proyecto desde una perspectiva empresarial.

Se llevaron a cabo reuniones con los tutores para entender los objetivos específicos del proyecto, como la detección y clasificación de acciones de movimiento (arriba, abajo, derecha, izquierda).

Se establecieron metas claras y medibles para el sistema de análisis de señales EEG, incluyendo los modelos a utilizar.

En esta fase se tomo la decisión de crear un modelo modular con un Notebook Principal que realizara llamadas al resto de Notebooks secundarios siguiendo los objetivos de negocio de poder ser escalables para futuras investigaciones.


\item
\textbf{Comprensión de los Datos}

Unos de los objetivos principales fue familiarizarme con los datos disponibles y realizar un análisis preliminar.

Hubo una exploración inicial donde se exploro los conjuntos de datos EEG para comprender su estructura, características y distribución.

Se hizo uso de herramientas y visualizaciones para identificar patrones y posibles anomalías en los datos.

Se identificaron valores atípicos (outliners) que podrían afectar negativamente el rendimiento de los modelos.

Se realizaron cursos online básicos propuestos en las reuniones de seguimiento:
Tutorial de Inicio de Pandas~\cite{curso:a}, Visualización de Datos~\cite{curso:b}, Trabajo con series temporales.~\cite{curso:c}

\item
\textbf{Preparación de los Datos}

El objetivo principal de esta fase es preprocesar y preparar los datos para su análisis en los diferentes modelos.

En esta fase se crearon nuevos conjuntos de datos para luego realizar análisis sobre ellos. En total 6 como se ha descrito con anterioridad.

Se aplicación de técnicas de normalización y escalado, como StandardScaler y Z-Score, para estandarizar las señales EEG.

Creación y utilización de ventanas deslizantes, agrupándolas por cada uno de los tipos de datos de cada conjunto de datos generado.

Al final de la fase se propuso la generación de Datos Sintéticos para aumentar la cantidad y diversidad del conjunto de datos y abordar el análisis de los modelos desde otra perspectiva.

\item
\textbf{Fase de Modelado}

El objetivo en esta fase fue seleccionar y aplicar técnicas de modelado adecuadas. Los que se han utilizado para el proyecto son los siguientes:

	\begin{itemize}
	
	\item
	\textbf{Algoritmos de Machine Learning:}

	K-Nearest Neighbors (KNN)
	Árboles de Decisión
	Random Forest

	\textbf{Evaluación y Comparación:}
	Uso de métricas como Tasa de acierto y loss para evaluar el rendimiento de estos modelos.

	\textbf{Redes Neuronales:}
	Multilayer Perceptron (MLP)
	Redes Neuronales Recurrentes (RNN)
	Long Short-Term Memory (LSTM)

	\textbf{Optimización y Regularización:} 

	Uso de técnicas como Dropout para mejorar la generalización de los modelos.
	\end{itemize}


Las dificultades que se encontraron en esta etapa fueron:

\textbf{- Compilación de modelos de redes neuronales. }

	Al ejecutar los primeros modelos de redes neuronales los datos normalizados me proporcionaban errores de compilación con modelos SRNN o LSTM, tuve que cambiar el escalado de los datos datos y shapear los modelos para que pudieran ejecutarse.

\textbf{- Utilización de ventanas temporales en los modelos de redes neuronales.} 

	En esta etapa no supe identificar este requerimiento por parte de los tutores y estuve implementando varias formas de poder utilizar ventanas temporales en el código, esto hizo que tuviera un gran retraso en la finalización y aceptación del codigo.

\textbf{- Utilización de modelos de redes neuronales y callbacks.} 

	El sobreajuste en los modelos predictivos era una constante y se implemento el uso de callbacks para evitar este sobreajuste.


\textbf{- Gráficas y definiciones básicas como normalizar o escalar el conjunto de datos.}



\item
\textbf{Evaluación}


En esta fase se trata de poder evaluar el modelo para asegurar que cumple con los objetivos del negocio. Para ellos se utilizaran las siguientes técnicas:


	\begin{itemize}
	
	\item
	\textbf{Métricas de Evaluación:}

Para todos los modelos se ha utilizado la evaluación mediante subconjuntos de datos de validación y prueba. (Val y Test)

	\item
	\textbf{Utilización de Callbacks y Técnicas de Monitorización:}

 Se ha hecho uso de EarlyStopping, ReduceLROnPlateau y ModelCheckpoint para evitar el sobreajuste y guardar el mejor modelo durante el entrenamiento.
	\end{itemize}


\item
\textbf{Despliegue}

En esta fase el objetivo es la implementacion del proyecto y para ello se empezó con la definición y creación de los primeros Jupyter Notebooks.


\textbf{Jupyter Notebooks.} 

Se han utilizado no solo para el desarrollo iterativo y la experimentación rápida, sino también para la implementación modular del proyecto. 

La capacidad de crear, documentar y ejecutar celdas de código de manera interactiva permite a los usuarios finales poder ajustar y probar el proyecto en tiempo real, facilitando la depuración y el ajuste fino.

La naturaleza modular de los Jupyter Notebooks ha permitido separar distintas fases del proceso, desde la carga y preprocesamiento de datos hasta el entrenamiento y evaluación del modelo. 

Esta separación ha facilitado el mantenimiento y la actualización de cada componente sin afectar al resto del sistema.

\end{enumerate}





\section{Estudio de viabilidad}


\subsection{Viabilidad económica}

La viabilidad económica evalúa la posibilidad de generar ingresos y recuperar la inversión realizada a través del proyecto desarrollado. Se han de considerar las siguientes áreas:
\begin{itemize}
\tightlist
   \item
    \textbf{Potencial del proyecto para Otros BCIs:} 
    
    El código desarrollado para el análisis de señales EEG podría adaptarse y aplicarse a una variedad de sistemas BCI disponibles en el mercado. La flexibilidad del lenguaje de programación Python permite su integración con sistemas existentes, ampliando su potencial de mercado.
    
   \item
    \textbf{Mercado Objetivo:} 
    
    Los fabricantes de sistemas BCI constituyen el mercado principal para la comercialización del software desarrollado. Ejemplos de fabricantes incluyen NeuroSky, Emotiv, Brain Products y otros proveedores de dispositivos de interfaz cerebro-ordenador utilizados en investigación, medicina y aplicaciones comerciales. En España destacan Starlab Barcelona, Bitbrain, Neuroelectrics y Starlab.
    
   \item
    \textbf{Modelo de Negocio:} 
    
    Se propone un modelo de negocio basado en licencias de software. Los fabricantes de BCI podrían adquirir licencias comerciales del software para integrarlo en sus productos o servicios. Las licencias podrían estructurarse en diferentes niveles de funcionalidad y soporte, permitiendo adaptarse a las necesidades específicas de cada fabricante.
    
   \item
    \textbf{Estrategia de Comercialización:}
    
    La estrategia de comercialización incluiría demostraciones técnicas, pruebas piloto y colaboraciones estratégicas con fabricantes de BCI. Se realizarían  participaciones y colaboraciones en conferencias y eventos de la industria, así como la utilización de plataformas online para promover el software y captar el interés de potenciales clientes.

\end{itemize}




\subsection{Viabilidad legal}


Tipos de licencias para el sofware y bibliotecas instalado:

\begin{itemize}
 \tightlist
  \item
   \textbf{BSD (Berkeley Software Distribution):} 
    
    BSD es una licencia permisiva de código abierto que permite el uso, modificación y distribución del software con pocas restricciones. Requiere que se mantengan los avisos de copyright y las cláusulas de la licencia..
    
  \item
   \textbf{ASF (Apache Software License):}
    
    ASF 2.0 permite el uso, modificación y distribución del software, otorgando derechos de patente explícitos. Requiere mantener los avisos de copyright y proporcionar una NOTIFICACIÓN de cambios significativos.
    
  \item
   \textbf{Código abierto (Open source):}
    
    Las licencias de código abierto permiten el acceso, uso, modificación y distribución libre del software. Estas licencias fomentan la colaboración y la transparencia en el desarrollo de software, garantizando la libertad de uso y mejora.
\end{itemize}



El estudio de viabilidad ha demostrado que el desarrollo del software de análisis de datos EEG presenta una sólida oportunidad económica y cumple con los requisitos legales necesarios para su comercialización.

La aplicación potencial del software a otros sistemas BCI y la flexibilidad del modelo de negocio basado en licencias ofrecen un camino claro hacia la monetización y la expansión del proyecto.

  A continuación se presenta una tabla con algunas de las principales bibliotecas y herramientas de software utilizadas en el proyecto, junto con el tipo de licencia y las restricciones asociadas:
 

\begin{table}[p]
	\centering
	\begin{tabularx}{\linewidth}{ p{0.3\columnwidth} p{0.3\columnwidth} p{\columnwidth}}
		\toprule
		\textbf{Software}    & \textbf{Licencia} & \textbf{Restricciones} \\
		\toprule		
		\textbf{Anaconda}              & Gratuita    & Uso personal.  \\
		\toprule	
		\textbf{Jupyter Notebook}              & Libre    & Código abierto \\
		\toprule			
		\textbf{Python}              & Libre permisiva    & BSD \\
		\toprule		
		\textbf{Numpy}              & Libre permisiva    & BSD \\
		\toprule		
		\textbf{Scipy}              & Libre permisiva    & BSD \\
		\toprule		
		\textbf{Pandas}               & Libre permisiva    & BSD \\
		\toprule		
		\textbf{Matplotlib}              & Libre permisiva    & BSD \\
		\toprule		
		\textbf{Seaborn}              & Libre permisiva    & BSD \\
		\toprule		
		\textbf{Sklearn}               & Libre permisiva    & BSD \\
		\toprule		
		\textbf{Tensorflow}              & Libre permisiva    & ASF \\
		\toprule		
		\textbf{Ipywidgets}              & Libre    & Código abierto \\
		\bottomrule
	\end{tabularx}
	\caption{Licencias Software}
\end{table}
\apendice{Especificación de Requisitos}

\section{Introducción}

La Especificación de Requisitos tiene como objetivo definir de manera clara y detallada las necesidades y expectativas del proyecto que consta de un notebook Jupyter definido en Python, el cual invoca a otros notebooks Jupyter secundarios para realizar diversas tareas. 

Este documento sirve para asegurar que todas las partes interesadas comprendan y acuerden los objetivos y funcionalidades del sistema a desarrollar.

En un entorno de trabajo basado en notebooks Jupyter, es esencial contar con una especificación precisa que guíe el desarrollo, implementación y validación del sistema. Esto no solo facilita la colaboración y la comunicación entre los desarrolladores, analistas de datos y usuarios finales, sino que también ayuda a identificar posibles riesgos y desafíos durante el ciclo de vida del proyecto.


\section{Objetivos generales}

Los objetivos principales de este TFG son los siguientes:

\begin{itemize}
\tightlist
\item
	\textbf{Facilitar el procesamiento de datos EEG.} Generados en la universidad mediante el interfaz BCI.  	
\item
 	\textbf{Automatizar tareas repetitivas.} Con la creacion de notebooks secundarios para poder ejecutar las tareas comunes y repetitivas relacionadas con los datos EGG. 
\item
	\textbf{Mejorar la tasa de acierto en los analisis de modelos} implementando modelos precisos y fiables. 
\item
	\textbf{Asegurar la accesibilidad y usabilidad en los notebooks.} Los notebooks son herramientas intuitivas y accessibles que aportan esa facilidad a la hora de interacturar.
\item
	\textbf{Integracion y extensibilidad.} Al utilizar llamadas a otros notebooks se asegura que se permita la integracion de nuevos notebooks y se asegura posibles nuevas mejoras en el codigo.
		
\end{itemize}



\section{Catálogo de requisitos}

Hay dos tipos de requisitos, los funcionales (qué debe hacer el código) y los no funcionales (cómo debe funcionar el código): 

\subsection{Requisitos funcionales}

\begin{itemize}
\tightlist
\item
  \textbf{RF-001 Iniciar notebook Main:}
 
  \begin{itemize}
  \tightlist
  \item
   \textbf{Descripción:} Los usuarios han de poder iniciar el notebook principal para poder comenzar con la ejecución del trabajo.
  \item
   \textbf{Prioridad:} Alta
  \item
   \textbf{Criterios de aceptación:} El notebook Main debe abrirse correctamente en menos de 5 segundos. 
  \end{itemize}


\item
  \textbf{RF-002 Invocación de Notebooks Secundarios Preparatorios:}
  
  \begin{itemize}
  \tightlist
  \item
    \textbf{Descripción:} Los usuarios han de poder invocar los notebook secundarios preparatorios para poder ejecutar con seguridad el resto de notebooks de preprocesado y modelados.
  \item
   \textbf{Prioridad:} Alta
  \item
   \textbf{Criterios de aceptación:} Los notebooks secundarios preparatorios se han de ejecutar correctamente y que la celda en la que se se llama al notebook secundario no devuelva errores. Los notebooks preparatorios. Su ejecucion no ha de ser mayor a 5 minutos.
  \end{itemize}
 
  
\item
  \textbf{RF-003 Invocación de Notebook Secundarios carga de archivo CSV:}

  \begin{itemize}
  \tightlist
  \item
    \textbf{Descripción:} Los usuarios han de poder invocar los notebook secundarios carga de archivo CSV, para la ejecucion de este notebook secundario se ha de subir un archivo csv con el conjunto de datos a analizar y una vez subido el archivo, se iluminara un botón con nombre 'Procesar datos' que procesara los datos una vez pulsado.
  \item
    \textbf{Prioridad:} Alta
  \item
    \textbf{Criterios de aceptación:} Si no se ejecuta correctamente la llamada a este notebook no se podrá seguir con la ejecución del resto de notebooks puesto que sin datos no se podrían realizar las siguientes acciones. 
  \end{itemize}


\item
  \textbf{RF-004 Invocación de Notebooks Secundarios análisis y preprocessing:}
  
  \begin{itemize}
  \tightlist
  \item
   \textbf{Descripción:} Los usuarios han de poder invocar los notebook secundarios análisis y preprocessing para poder proporcionar este procesado de datos antes de la ejecución de modelos machine o deep learning.
  \item  
   \textbf{Prioridad:} Alta
  \item  
   \textbf{Criterios de aceptación:} Si no se ejecuta correctamente la llamada a este notebook no se podrá seguir con la ejecución del resto de notebooks puesto que sin la realizacion del preprocessing de los datos no se generan los csv correcpondientes para poder seguir ejecutando los siguientes notebooks.
  \end{itemize}


\item
  \textbf{RF-005 Invocación de Notebooks Secundarios para experimentos:}

  \begin{itemize}
  \tightlist
  \item
   \textbf{Descripción:} Los usuarios han de poder invocar los notebook secundarios para experimentos y poder generar los datos sobre Tasa de acierto y generacion de matrices de confusión.
  \item  
   \textbf{Prioridad: Media}
  \item  
   \textbf{Criterios de aceptación:} Estos notebooks se puede ejecutar de manera individual o en el orden pre establecido. para realizar el análisis final lo ideal es ejecutar todos los notebooks de esta session.
	\textbf{}
  \end{itemize}

\item
  \textbf{RF-006 Invocación de Notebook Secundario resultados:}

  \begin{itemize}
  \tightlist
  \item
   \textbf{Descripción:} Los usuarios han de poder invocar los notebook Secundario resultados y poder visualizar los resultados de las particiones test del conjunto de datos. Se han de haber ejecutado todos los notebooks de esta sección.
  \item  
   \textbf{Prioridad: Alta}
  \item  
   \textbf{Criterios de aceptación:} Sin la ejecución de este notebook no se obtendrian los resultados del trabajo y por lo tanto es critico este ultimo notebook para poder visualizar la tasa de acierto en datos que el modelo no conocía en la etapa de entrenamiento y validación.
  \end{itemize}


\end{itemize}


\subsection{Requisitos no funcionales}

\begin{itemize}
\tightlist
\item
  \textbf{RNF-001 Rendimiento:}
 
  \begin{itemize}
  \tightlist
  \item
   \textbf{Descripción:} El sistema donde se ejecuten los notebooks ha de poder manejar la ejecución sin causar demoras significativas.
  \item
   \textbf{Criterios de aceptación:} El tiempo de ejecución para la ejecución de cada notebook secundario no ha de ser superior a 15 minutos.
  \end{itemize}


\item
  \textbf{RNF-002 Usabilidad:}
  
  \begin{itemize}
  \tightlist
  \item
    \textbf{Descripción:} Los notebooks han de ser faciles de entender y usar por los usuarios.
  \item
   \textbf{Criterios de aceptación:} La interfaz de los notebooks ha de ser intuitiva y clara.
  \end{itemize}
 
  
\item
  \textbf{RNF-003 Escalabilidad:}

  \begin{itemize}
  \tightlist
  \item
    \textbf{Descripción:} En la configuración del notebooks principal se debe permitir la adhesión de nuevos notebooks y sus llamadas desde el notebook principal.
  \item
  \item
    \textbf{Criterios de aceptación:} Poder seguir integrando nuevos notebooks para ser llamados desde el notebook principal sin afectar al funcionamiento del código. 
  \end{itemize}

\end{itemize}



\section{Especificación de requisitos}


% Caso de Uso 1 -> Consultar Experimentos.
\begin{table}[p]
	\centering
	\begin{tabularx}{\linewidth}{ p{0.21\columnwidth} p{0.71\columnwidth} }
		\toprule
		\textbf{CU-1}    & \textbf{Ejemplo de caso de uso}\\
		\toprule
		\textbf{Versión}              & 1.0    \\
		\textbf{Autor}                & José Luis Pérez Gómez \\
		\textbf{Requisitos asociados} & RF-xx, RF-xx \\
		\textbf{Descripción}          & La descripción del CU \\
		\textbf{Precondición}         & Precondiciones (podría haber más de una) \\
		\textbf{Acciones}             &
		\begin{enumerate}
			\def\labelenumi{\arabic{enumi}.}
			\tightlist
			\item Pasos del CU
			\item Pasos del CU (añadir tantos como sean necesarios)
		\end{enumerate}\\
		\textbf{Postcondición}        & Postcondiciones (podría haber más de una) \\
		\textbf{Excepciones}          & Excepciones \\
		\textbf{Importancia}          & Alta o Media o Baja... \\
		\bottomrule
	\end{tabularx}
	\caption{CU-1 Nombre del caso de uso.}
\end{table}



\apendice{Especificación de diseño}

\section{Introducción}

La especificación de datos es un componente crítico en el desarrollo de sistemas de información, especialmente cuando se trabaja con conjuntos de datos complejos como los datos EEG. 

Mediante BCI, sistema que permite mediante adquisición de señales EEG, se han podido interpretar las señales adquiridas a través de las señales EEG y poder transformarlas en un conjunto de datos para su posterior análisis.  



\textbf{Descripción de los Datos:}

El archivo datosEEGTotal.csv contiene los datos facilitados para poder realizar los experimentos para la ejecución del TFG.

Su formato es de tipo CSV con un separador (;) entre los datos que lo componen.

Las características recogidas en el archivo de los datos EEG son las siguientes:

Timestamp, Attention, Meditation, Delta, Theta, LowAlpha, HighAlpha, LowBeta, HighBeta, LowGamma, HighGamma, Signal y Key.

\textbf{Timestamp:} Registro de tiempo para los experimentos, medido en mili-segundos.

\textbf{Attention:} Registra el grado de atención del participante que realiza el experimento.

\textbf{Meditation:} Grado de calma que tendría el individuo.

\textbf{Delta:} Son ondas de baja frecuencia (1 y 4 Hz), están presentes en etapas de sueño profundo, durante una meditación profunda y en pacientes con lesiones cerebrales o con TDAH severo.

\textbf{Theta:} Ondas entre 4 y 8 Hz, se encuentran en estados de calma profunda y sueño R.E.M., están ligadas al aprendizaje, memoria y intuición.

\textbf{Alpha:} Ondas entre 8 y 12 Hz, representan un estado de poca actividad cerebral y se asocian a un estado de calma mental. Divididas en dos señales LowAlpha y HighAplpha

\textbf{Beta:} Se diferencian en LowBeta y HighBeta, su frecuencia esta entre 12 y 35Hz, asociadas a una alta actividad mental. 

\textbf{Gamma:} En los datos se diferencia LowGamma y HighGamma, son ondas por encima de 30Hz y suelen aparecer cuando hay una alta concentración o atención

\textbf{Signal:} Podría ser la señal de que aporta la interfaz BCI.

\textbf{Key:} Valores target de lo que el individuo estaba pensando o visualizando durante el experimento.


La transformación de datos o el preprocessing que se ha realizado para poder afrontar los experimentos del TFG han sido los siguientes:


\textbf{Unificación de características Key:} El conjunto de datos tiene varios valores en Key que indican los mismo. LButton y Left.

\imagen{anexos/Key}{Cantidad de elementos para la característica Key}

\textbf{División del conjunto de datos:} El conjunto de datos tiene cuatro segmentos divididos por su Timestamp, superponiéndose entre ellos. Divido en estos cuatro segmentos para poder realizar experimentos y también dejo el conjunto de datos sin dividir para realizar experimentos conjuntos a los cuatro segmentos.

\imagen{anexos/Timestamp}{Análisis característica Timestamp del conjunto de datos}

\textbf{Eliminación de características:} Elimino las características Signal por no aportar nada significativo en el conjunto de datos y Timestamp porque no quiero que los datos aportados puedan tener una patrón temporal que haga que los experimentos no sean reales.


\textbf{Escalado de datos:} He utilizado la opción de escalar los datos ya que no son datos normales puesto que no tienen una distribución gaussiana en sus datos. 

\section{Diseño de datos}




\section{Diseño procedimental}



\section{Diseño arquitectónico}


\apendice{Documentación técnica de programación}

\section{Introducción}

Este TFG de análisis de un conjunto de datos EEG se centra en la aplicación de técnicas avanzadas de aprendizaje automático utilizando notebooks Python. El objetivo principal es procesar y analizar datos EEG para identificar patrones y correlaciones significativas que puedan revelar comportamientos específicos en la actividad cerebral a la hora de realizar un experimento el individuo.

El propósito fundamental es aplicar modelos de aprendizaje automático para mejorar la comprensión de las señales EEG, los datos aportados por la universidad de burgos constan de unos experimentos realizados con BCI y varios voluntarios es los que estos usuarios deben visualizar una serie de imagenes compuestas por direcciones, arriba, abajo, derecha e izquierda.

Este trabajo podría tener aplicaciones importantes en tecnologías de interfaz cerebro-computadora (BCI) aplicadas a por ejemplo,  control de sillas de ruedas, control de dispositivos electrónicos, realidad virtual, videojuegos, control de robots, etc...


\section{Estructura de directorios}

La estructura de los directorios del trabajo es la siguiente:

  \begin{itemize}
  \tightlist
  \item
   \textbf{/documentacion/:} Esta carpeta contiene toda la documentación relacionada con el TFG.
  \item
   \textbf{/documentacion/imagenes/:} Archivo de imágenes para todo el proyecto, documentación y código.
  \item
   \textbf{/documentacion/latex:} Documentación en formato latex y PDF.
  \item
   \textbf{/codigo:} Carpeta para archivar notebooks y datos para la ejecución del código.
  \item
   \textbf{/codigo/datos/:} Archivo datosEEGTotal.csv con el origen de datos.
  \item
   \textbf{/codigo/datos/csv/:} Archivos de datos csv auto generados tras las ejecuciones de algunos notebooks.
  \item
   \textbf{/codigo/notebooks:} Contiene los notebooks Jupyter y por lo tanto, el código.
  \end{itemize} 
  
  
  

\section{Manual del programador}

Este manual sera utilizado principalmente por las personas involucradas en futuros cambios, mejoras o nuevos desarrollos en el código.

Para ello, detallo las funciones de cada uno de los componentes :


\section{Compilación, instalación y ejecución del proyecto}

\section{Pruebas del sistema}

\apendice{Documentación de usuario}

\section{Introducción}

\section{Requisitos de usuarios}

\section{Instalación}

\section{Manual del usuario}



\apendice{Anexo de sostenibilización curricular}

\section{Introducción}
Este anexo incluirá una reflexión personal del alumnado sobre los aspectos de la sostenibilidad que se abordan en el trabajo.
Se pueden incluir tantas subsecciones como sean necesarias con la intención de explicar las competencias de sostenibilidad adquiridas durante el alumnado y aplicadas al Trabajo de Fin de Grado.

Más información en el documento de la CRUE \url{https://www.crue.org/wp-content/uploads/2020/02/Directrices_Sosteniblidad_Crue2012.pdf}.

Este anexo tendrá una extensión comprendida entre 600 y 800 palabras.



\bibliographystyle{plain}
\bibliography{bibliografiaAnexos}

\end{document}
